\documentclass{IMW_poster}
% \documentclass{article}

% General packages
\usepackage[utf8]{inputenc}
\usepackage[english]{babel}
\usepackage{listings}
\usepackage{xcolor}
\usepackage[linktocpage=true, hidelinks]{hyperref}

% Math and Code Packages
\usepackage{amsfonts}
\usepackage{amssymb}
\usepackage{amsmath}
\usepackage{bm}
\usepackage{csquotes}
\usepackage{mathtools}
\usepackage{listings}
\usepackage{dsfont}
\usepackage[linesnumbered,ruled,vlined,algosection]{algorithm2e}

% Citation Packages
\usepackage[backend=biber,sorting=nty,natbib=true,style=authoryear]{biblatex}
\addbibresource{IMW_poster_template.bib}

% Definitions of Custom Math Commands
\DeclareMathOperator{\eps}{\varepsilon}
\DeclareMathOperator{\R}{\mathbb{R}}
\DeclareMathOperator{\E}{\mathbb{E}}
\DeclareMathOperator{\Prob}{\mathbb{P}}

% Information on the poster
\title{Example Posters: Working with the IMW Poster Template using the DOT and BiGSEM Logo}
\author{
    \textbf{First Author}, University of Vienna, 
    \texttt{first.author@univie.ac.at}\\
    \textbf{Second Author}, University of Vienna, 
    \texttt{second.author@univie.ac.at}
}
\abstract{%
    This is the abstract. To place the abstract on the poster, write the 
    text in the \texttt{\textbackslash abstract} command like in the 
    accompanying tex-file.\\
    This is a very, very long paragraph to show that the abstract spans over
    all columns and does not conform to the separation like the main body of 
    the poster. In general, the abstract should contain a concise description 
    and/or summary of the contents of the poster.\\
    This is a very, very long paragraph to show that the abstract spans over
    all columns and does not conform to the separation like the main body of 
    the poster. In general, the abstract should contain a concise description 
    and/or summary of the contents of the poster. 
}
\bottomright{
    \includegraphics[width=0.05\paperwidth]{example-image-a}
}


\begin{document}
    


\section{Basic Options}
Options are set in brackets with the \texttt{documentclass} command, like
\texttt{\textbackslash documentclass[{\normalfont\emph{<options>}}]\{IMW\_poster\}}.

\subsection{Paper Size}
Available papersizes are \texttt{a0} through \texttt{a4}, where \texttt{a4} is 
a draft-like option scaling the contents down. The default option is 
\texttt{a0}.

\subsection{Font Sizes}
The available font sizes are \texttt{14pt, 17pt, 20pt, 25pt} and 
\texttt{30pt}. The default option is \texttt{25pt}.

\subsection{Recommended Combinations}
The recommended combinations of paper size and font size are 
\texttt{\{a0, 25pt\}, \{a1, 17pt\}, \{a2, 14pt\}, \{a3, 14pt\},
\{a4, 25pt\}}. A4 papersize should be paired with a large font like
\texttt{30pt} or \texttt{25pt} since the contents are scaled down.\\
Note that the recommended combinations do not scale font size and paper size
perfectly.

\subsection{Logos in the Header}
The logos in the header of this poster template work similar to the logos in
a typical beamer template.
Each logo, from left to right \texttt{department}, \texttt{institute} and
\texttt{organization}, has a language option. The logos are set-up in such a
way that the whitespaces are balanced.
The \texttt{department} option currently allows for the \texttt{univie} and
\texttt{none} keys. The \texttt{institute} and \texttt{organization} options
currently default to \texttt{none}.\\
The default options are the German Uni Wien department logo and no institute
or organization logo.\\
Section \ref{sec:implementation} contains notes on how to add additional logos.

\subsection{The \texttt{structurecolor} option}
The \texttt{structurecolor} option can be called with any named color or color
mix as the key. It changes the color of structural elements like the section or
title headers and the line around main body.

\subsection{The \texttt{BiGSEM} and \texttt{CUDE} options}
The \texttt{BiGSEM} and \texttt{CUDE} options are aggregate options that set
the appropriate oragnization logos and set the structure color as the primary
color of the chosen logo.

\subsection{The \texttt{bottomright} option}
If the \texttt{bottomright} option is called, whatever is given in the 
\texttt{\textbackslash bottomright} command is placed in a \texttt{TikZ} node
at the bottom right. This is analogous to how the old poster template lookes.


\section{Figures and Tables}
To place a picture or table on the poster use the \texttt{figure} or 
\texttt{table} environments, respectively, as shown below. Specifying any 
placement options, will cause the figure or table to either not appear or 
display at the end of the document.

Following is an example figure and an example table:
\begin{figure}
    \centering
    \includegraphics[height=0.1\textheight]{example-image-a}
    \caption{This is an example figure.}
    \label{fig:example-figure}
\end{figure}

\setlength{\tabcolsep}{2ex}
\renewcommand{\arraystretch}{1.25}
\begin{table}
    \centering
    \begin{tabular}{p{0.35\linewidth}|p{0.2\linewidth}p{0.225\linewidth}}
        Name of Director & Time of Directorate & Num of Nobel Prizes (so
        far) \\
        \hline
        C. C. von Weizsäcker & 1970 - 1973 & 0 \\
        R. Selten & 1973 - 1984 & 1\\
        J. Rosenmüller & 1984 - 2002 & 0\\
        W. Trockel & 2002 - 2009 & 0\\
        F. Riedel & since 2009 & 0\\
    \end{tabular}
    \caption{This is an example table.}
    \label{tab:example-table}
\end{table}


\section{Math Environments}
Math environments work like usual:\\
For all $x \in \R$ and $Y \sim \mathcal{N}(0, 1)$ we have
\begin{align*}
    - x^2   &\leq 0 \\
            &\leq \Prob \left( Y \geq 0 \right) \leq 1.
\end{align*}
\begin{equation*}
    \sum_{i=1}^{\infty} \frac{1}{2^i} = 1
\end{equation*}
\begin{equation}\label{eq:example-equation}
    \exp\left( i \pi \right) + 1 = 0
\end{equation}


\section{References and Citations}
References work like ususal:\\
Tables like Table \ref{tab:example-table}, figures like Figures 
\ref{fig:example-figure} and equations like Equation 
\eqref{eq:example-equation} can be referenced.\\
Citations with \texttt{biblatex} can also be implemented: \cite{einstein},
\cite{dirac}, \cite{knuth-acp}, \cite{knuth-fa}, \cite{knuthwebsite}, 
\cite{ctan}


\section{Notes on Implementation}\label{sec:implementation}
\textbf{Logos in Header:} Adding additional department, institute or 
organization logos is simple. The process for adding a new institute is 
described below but departments or organizations are analogous. (Note that 
no department language is equivalent having just the base logo and thus the 
none language variant can be omitted.) Say we want to add the Example 
Institute and have logos as .png files on hand. In this case, we need to 
create a new folder named 'example' inside of logos/institutes, in which we 
put the logos and rename them to \texttt{example\_logo\_\{language\}.png}\\
For the logo without any text we replace \{language\} with 'none'.
Now we can simply call the theme with the option \texttt{institute=example}
and the desired `institutelang' like the pre-configured institutes.

\textbf{Floats:} Since the poster template is realized with a
\texttt{multicols} environment, floats are not supported in the main body of
the poster. As a result, specifying any placement options for figures or 
tables, will cause the figure or table to either not appear or display at the 
end of the document.


\printbibliography

\end{document}